\subsection{معماری و پیاده‌سازی}
پیاده‌سازی این الگوریتم در مسیر \lr{neuralops/compression\_engine/rpca\_compressor.py} قرار داشته و با نام‌های \lr{rpca} و \lr{rpca\_compressor} در رجیستری فشرده‌سازی فراخوانی می‌شود. از ویژگی‌های بارز معماری و مهندسی این ماژول می‌توان به موارد ذیل اشاره کرد:
\begin{itemize}
    \item \textbf{محاسبات پایدار با دقت \lr{FP32}:} تمامی مراحل چرخه‌های تکرار \lr{PCP} و تجزیه‌های \lr{torch.linalg.svd} در فرمت ممیز شناور ۳۲ بیتی محاسبه شده و مقادیر منفرد با کف $10^{-12}$ جهت ممانعت از ناپایداری عددی محدود می‌شوند.
    \item \textbf{سازگاری کامل با لایه‌های خطی و کانوولوشن ۱ بعدی (\lr{Conv1D}):} تشخیص خودکار لایه‌های کانوولوشن ۱ بعدی (مانند ساختار \lr{GPT-2}) از طریق وارسی نام کلاس و ویژگی \lr{nf}، ترانهش صحیح ماتریس‌ها و نگاشت ابعاد را تضمین می‌کند.
    \item \textbf{مکانیزم تنزیل اضطراری دومرحله‌ای (\lr{Robust Fallback Hierarchy}):} در صورت بروز خطا در چرخه‌های \lr{ALM}، الگوریتم ابتدا به \lr{SVD} مستقیم روی ماتریس خام رجوع کرده و در صورت عدم دستیابی به شرط پایداری، لایه اصلی را بدون تخریب حفظ می‌نماید.
    \item \textbf{پایش تفکیکی خطای فروبنیوس:} خطای بازسازی نسبی فروبنیوس $\epsilon_{\text{recon}} = \frac{\|\mathbf{W} - \mathbf{B} \mathbf{A}\|_F}{\|\mathbf{W}\|_F}$ برای تک‌تک لایه‌ها ثبت و پایش می‌شود.
\end{itemize}
